\documentclass{ctexbook}
\usepackage{ctex}
\usepackage{amsmath}
\usepackage{amssymb}
\usepackage[colorlinks,linkcolor=blue]{hyperref}
\usepackage{booktabs}
\usepackage[paper=a4paper, left=1.5cm, right=1.5cm, top=3cm, bottom=3cm]{geometry}
\usepackage{enumitem}
\usepackage{tikz}

\newcommand{\no}{\noindent}
\newcommand{\pro}{\no{\bf 证明：}}


\newcommand{\ITEMS}{\begin{enumerate}[leftmargin=4em]}
\newcommand{\ITEME}{\end{enumerate}}


\title{实变函数论}

\begin{document}

\maketitle

\tableofcontents

\chapter{预备知识}

尽管黎曼积分很有用，但也存在很大的局限，我们旨在建立一种新的积分，其应用范围更广，并希望它可以看作黎曼积分的推广。

% 具体描述黎曼积分的缺陷

\section{预备知识：数学分析}

\subsection{数列的上、下极限}

\subsubsection{上、下确界的定义}

粗略地讲，对于$A \subseteq \mathbb{R}$，$A$ 的上确界即为$A$的最小上界,记作
\[
\sup A
\]
$A$的下确界为$A$的最大下界，记作
\[
\inf A
\]
\smallskip

% 数学分析九大定理之一，作为定理使用
\subsubsection{单调有界序列必收敛}
设$\{a_n \}_{n\ge 1}$为单调数列，现在对其单调性进行分类讨论：
\ITEMS
	\item 当$\{a_n \}_{n\ge 1}$单调递增时：若数列有界，则数列极限为上确界，即$\lim\limits_{n\to\infty} a_n = \sup\{a_n\}$；若数列无界，数列极限不存在，即$\lim\limits_{n\to\infty} a_n = +\infty$。
	\item 当$\{a_n \}_{n\ge 1}$单调递减时：若数列有界，则数列极限为下确界，即$\lim\limits_{n\to\infty} a_n = \inf\{a_n\}$；若数列无界，数列极限不存在，即$\lim\limits_{n\to\infty} a_n = -\infty$。
\ITEME

\smallskip

\subsubsection{上、下极限的定义}

\no {\bf 定义1}

设$\{a_n \}_{n\ge 1}$为一数列，其所有极限点（收敛子列的极限） 的上确界定义为上极限，记作
\[
	\varlimsup _{n\rightarrow \infty }a_{n}
\]
同样的，所有极限点的下确界为下极限，记作
\[
	\varliminf _{n\rightarrow \infty }a_{n}
\]
\smallskip

\no {\bf 定义2}

构造$\{b_m\}_{m\ge1}$，令$b_m = \sup\{a_m,a_{m+1},\cdots\}$，显然可知$\{b_m\}$单调递减，原数列的上极限可以定义为
\[
\begin{aligned}
	\varlimsup _{n\rightarrow \infty }a_{n} &= \lim_{m\to\infty} b_m = \inf\{b_m\}\\
	&= \lim_{m\to\infty}\left(\sup_{n\ge m}{\{a_n\}}\right)\\
	&= \inf_{m\ge1}\left(\sup_{n\ge m}{\{a_n\}}\right)\\
	&= \limsup _{n\rightarrow \infty }\{a_{n}\}
\end{aligned}
\]
类似的，我们可以构建$\{c_m\}_{m\ge1}$，令$c_m = \inf\{a_m,a_{m+1},\cdots\}$，显然可知$\{c_m\}$单调递增，原数列的下极限可以定义为
\[
\begin{aligned}
	\varliminf _{n\rightarrow \infty }a_{n} &= \lim_{m\to\infty} c_m = \sup\{c_m\}\\
	&= \lim_{m\to\infty}\left(\inf_{n\ge m}{\{a_n\}}\right)\\
	&= \sup_{m\ge1}\left(\inf_{n\ge m}{\{a_n\}}\right)\\
	&= \liminf _{n\rightarrow \infty }\{a_{n}\}
\end{aligned}
\]

\subsubsection{上、下极限的性质}

\no{\bf 性质1}

设$\{d_k\}$是$\{a_n\}$的收敛子序列，则有
\[
	\varliminf _{n\rightarrow \infty}a_{n} \le \lim_{k\to\infty}d_k \le \varlimsup _{n\rightarrow \infty}a_{n}
\]

\no{\bf 性质2}

数列$\{a_n\}$收敛的充要条件是，数列的上下极限相等，且为有限实数，即
\[
	\lim_{n\to\infty}a_n = L \iff \varliminf _{n\rightarrow \infty}a_{n} = \varlimsup _{n\rightarrow \infty}a_{n} = L \in \mathbb{R}
\]

\subsection{函数列的上、下极限}

对于在区间$I$上的函数列，上极限是一个在区间$I$上的函数，记作
\[
	\varlimsup_{n\to\infty} f_n(x)
\]
而其函数取值如下，其中$\left\{f_1(x),f_2(x),\cdots\right\}$是将自变量$x$带入函数列得到的数列
\[
	\varlimsup_{n\to\infty} f_n(x) = \varlimsup\left\{f_1(x),f_2(x),\cdots\right\}
\]
同理我们也可以得到如下的下极限
\[
	\varliminf_{n\to\infty} f_n(x) = \varliminf\left\{f_1(x),f_2(x),\cdots\right\}
\]
\smallskip

与数列极限的性质相似，我们可以得到
\[
	\varliminf_{n\to\infty} f_n(x) \le \varlimsup_{n\to\infty} f_n(x)
\]
\bigskip

\section{预备知识：集合论}

\subsection{集合基本概念}

\subsubsection{基本集合}

% 在课上的全集使用的是X 而空集是\Phi
我们将用$X$表示全集，$\emptyset$表示空集

\subsubsection{集合簇}

我们称由集合组成的集合为集合簇，记作
\[
	\{A_\alpha\}_{\alpha\in I}
\]
其中$A_\alpha \subseteq X$，$I$是指标集（指标集不一定是$\mathbb{N}_+$也可以是$\mathbb{R}$，或其他任意集合）

\subsection{集合运算}

\subsubsection{基础运算}

交集记作$A\cap B$，并集记作$A\cup B$，补集记作$A^c$

\subsubsection{减法运算}

集合减法记作$A\setminus B = A\cap B^c$

\subsubsection{对称差}

对称差记作$A\triangle B = A\oplus B = (A\setminus B)\cup(B\setminus A)$

\subsubsection{大型运算}

\[
\begin{aligned}
	\bigcup_{n=1}^\infty\left(a+\frac1n,\infty\right)&=(a,+\infty)\\
	\bigcap_{n=1}^\infty\left(-\infty,a+\frac1n\right]&=(-\infty,a]\\
	\bigcap_{n=1}^\infty\left(a-\frac1n,\infty\right)&=[a,+\infty)\\
	\bigcup_{n=1}^\infty\left(-\infty,a-\frac1n\right]&=(-\infty,a)
\end{aligned}
\]

\subsection{集合列的上、下极限}

设$\{A_n\}_{n\ge1}$为集合列，其中$\forall A_n\subseteq X$。为定义集合列的上极限，我们构造集合列$\{B_m\}$如下：
\[
	B_m = \sup_{n\ge m}\{A_n\} = \bigcup_{n\ge m} A_n
\]
按包含关系，显然$B_1\ge B_2\ge \cdots$，可定义
\[
	\varlimsup_{n\to\infty}\{A_n\}=\lim_{m\ge 1}\{B_m\}=\bigcap_{m\ge1}B_m=\bigcap_{m\ge1}\bigcup_{n\ge m} A_n
\]
类似的，我们可以构造集合列$\{C_m\}$ 如下：
\[
C_m = \inf_{n\ge m}\{A_n\} = \bigcap_{n\ge m} A_n
\]
按包含关系，显然$C_1\le C_2\le \cdots$，可定义
\[
	\varliminf_{n\to\infty}\{A_n\}=\lim_{m\ge 1}\{C_m\}=\bigcup_{m\ge1}C_m=\bigcup_{m\ge1}\bigcap_{n\ge m} A_n
\]
我们也可以将上、下极限分别定义为
\[
	\varlimsup_{n\to\infty}\{A_n\}=\{x\in X: x\in \mbox{无穷多个}A_n\}
\]
\[
	\varliminf_{n\to\infty}\{A_n\}=\{x\in X: \exists\,N\in\mathbb{N}\ s.t.\ x\in A_n,\forall n\ge N\}
\]
这里我们只给出两种下极限定义等价性的证明（上极限的情况留给大家作思考题）: 
\[ 
    x\in\varliminf_{n\to\infty}\{A_n\}=\bigcup_{m\ge1}C_m\Leftrightarrow \exists\,N\in\mathbb{N}\ s.t.\ x\in C_N=\bigcap_{n\ge N} A_n \Leftrightarrow x\in A_n,\forall n\ge N
\]
显然
\[
	\varliminf_{n\to\infty}\{A_n\} \subseteq \varlimsup_{n\to\infty}\{A_n\}
\]
\subsection{映射}
设$X$、$Y$是两个集合，如果$X$中任一个元素$x$都按照某种规律$f$对应到$Y$的某一元素（记为$f(x)$），我们则称$f$为一从$X$到$Y$的映射，通常记为：
\[
	f:X\rightarrow Y
\]
对于$\forall A\subseteq X$，我们称$f(A)=\{f(x):x\in A\} $为$A$的像；对于$\forall B\subseteq Y$，我们称$f^{-1}(B)=\{x\in X:f(x)\in B\} $为$B$的原像
\subsubsection{单射}
满足以下条件的映射称为单射
\[
	x_1 \ne x_2 \Rightarrow f(x_1) \ne f(x_2)
\]
\subsubsection{满射}
满足以下条件的映射称为满射
\[
	f(X)=Y
\]
\subsubsection{双射}

既是单射又是满射的映射称为双射（或一一对应）
\subsubsection{逆映射}

当映射$f:X\rightarrow Y$是双射时，$f^{-1}$ 成为映射，称为 $f$ 的逆映射，记为
\[
	f^{-1}:Y\rightarrow X
\]

\subsubsection{映射与集合运算}

对于任何映射$f:X\rightarrow Y$，有集合列$\{A_\alpha\subseteq X\}_{\alpha\in I}$，$\{B_\alpha\subseteq X\}_{\alpha\in I}$，我们有
\[f\left(\bigcup_{\alpha\in I} A_\alpha\right) = \bigcup_{\alpha\in I}f(A_\alpha)\]
\[f\left(\bigcap_{\alpha\in I} A_\alpha\right) \subseteq \bigcap_{\alpha\in I}f(A_\alpha)\]
\[f^{-1}\left(\bigcup_{\alpha\in I} A_\alpha\right) = \bigcup_{\alpha\in I}f^{-1}(A_\alpha)\]
\[f^{-1}\left(\bigcap_{\alpha\in I} A_\alpha\right) = \bigcap_{\alpha\in I}f^{-1}(A_\alpha)\]
\[f^{-1}\left(B_\alpha^c\right) = \left(f^{-1}(B_\alpha)\right)^c\]


\subsection{集合的基数（势）}

粗略地讲，任一集合 $A$ 的基数或者势，记为 $|A|$，是指 $A$ 所含元素的个数。我们称自然数集
\[
	\mathbb{N}=\{1,2,\cdots\}
\]
的基数为可数无穷，记为 $\aleph_0$（读作：阿列夫0）。

在比较任意两个集合 $A$ 和 $B$ 的基数时，我们有下面的定义：

\no {\bf 定义1} 我们称 $\lvert A \rvert \le \lvert B \rvert$，如果存在从 $A$ 到 $B$ 的单射。\\
\no {\bf 定义2} 我们称 $\lvert A \rvert = \lvert B \rvert$，如果$\lvert A \rvert \le \lvert B \rvert$ 且 $\lvert B \rvert \le \lvert A \rvert$，也就是说，存在从 $A$ 到 $B$ 的双射。\\

\no {\bf 例} 记 $\mathbb{N}_2=\{2,4,\cdots\}$，即所有偶自然数的集合。不难看出 $\lvert\mathbb{N}_2\rvert=\lvert\mathbb{N}\rvert=\aleph_0$，因为每个自然数乘 $2$ 就是一从 $\mathbb{N}$ 到 $\mathbb{N}_2$ 的双射。

\subsubsection{幂集}
称任一集合$A$的所有子集构成的集合为集合$A$的幂集，记为$P(A)$。不难看出
\[
\lvert P(A) \rvert = 2^{\lvert A \rvert}
\]

\subsubsection{无最大基数定理}

对于任一非空集合 $A$，我们总有 $\lvert A \rvert<\lvert P(A) \rvert$。
\smallskip

\pro

我们先证明 $\lvert A \rvert \le \lvert P(A) \rvert$：\\
定义映射 $f:A\rightarrow P(A)$ 如下：
\[ 
	g(x)=\{x\}, \forall x \in A
\] 
不难看出，$g$ 是单射。

下面我们用反证法证明 $\lvert A \rvert \ne \lvert P(A) \rvert$:
若相等，则存在双射
\[
	f:A\rightarrow P(A)
\]
定义一个$A$的子集合$B$如下：
\[
B = \{x\in A: x\notin f(x)\}
\]
由于 $f$ 是双射，所以必存在（唯一）元素 $y\in A$ 使得
\[
f(y) = B
\]
考虑 $y$ 与 $B$ 的关系，发现无论如何都会产生矛盾：
\ITEMS
	\item 如果 $y\in B=f(y)$, 由 $B$ 的定义，我们会得到 $y\notin B$；
	\item 如果 $y\notin B=f(y)$, 由 $B$ 的定义，我们会得到 $y\in B$。
\ITEME

\no {\bf 注} 我们记 $\aleph_n=2^{\aleph_{n-1}}, n=1, 2, \cdots$。由上述定理，显然我们有 
\[
	\aleph_0 < \aleph_1 < \aleph_2 <\cdots< 
\]

\subsubsection{关于基数的几个基本定理}
\ITEMS
\item 对任意集合 $A$ 和 $B$，若 $A\subseteq B$，则 $\lvert A \rvert \le \lvert B \rvert$。
\item 设 $A$ 是可数集（有限集或可数无穷集），$B$ 是无穷集，则 $\lvert A\cup B \rvert= \lvert B \rvert$。
     （例：Hilbert旅馆）
\item 设集合列 $\{A_n\}_{n=1}^{\infty}$ 中每一 $A_n$ 的基数都是 $\aleph_0$，则 $\lvert \bigcup_{n=1}^\infty A_n \rvert=\aleph_0 $。也就是说，可数个可数集的并集还是可数集。 （例：记 $\mathbb{Q}=$ 有理数集，则 $\lvert\mathbb{Q}\rvert=\aleph_0 $） 
\ITEME


\subsubsection{实数集的基数}
实数集 $\mathbb{R}$ 的基数一般记为 $c$。 证明 $c=\aleph_1$。\\
\pro

我们先证明 $\lvert \mathbb{R} \rvert=\lvert (0,1]\rvert $，分三步进行：\\

1. $\lvert \mathbb{R}\rvert =\lvert (-1,1)\rvert $: 不难看出映射  $f:(-1,1)\rightarrow \mathbb{R}$ 满足 $f(x)=\frac {x}{1-x^2}$ 是双射。\\

2. $\lvert (-1,1)\rvert=\lvert (0,1)\rvert $: 不难看出映射  $f:(-1,1)\rightarrow (0,1) $ 满足 $f(x)=\frac {x+1}{2}$ 是双射。\\

3. $\lvert (0,1)\rvert=\lvert (0,1]\rvert $： 因为 $(0,1)\subseteq (0,1]\subseteq (-1,1)$。\\

下面我们证明 $\lvert (0,1]\rvert=2^{\aleph_0}$: \\
1. 把 $(0,1]$ 中的任一实数 $x$ 用二进制表示：更具体地讲，就是找到每项取值是 $0$ 或 $1$ 的无穷数列 $\{a_n\}_{n \ge 1}$ 满足在二进制下 $x=0.a_1 a_2\cdots a_n \cdots$，其中无穷多个 $a_n$ 是 $1$ 。\\
\ITEMS 
	\item 第一步: 把区间 $(0,1]$ 二等分为两半开闭子区间：若 $x$ 属于靠左的（称为第一个，后同）子区间 $(0,\frac{1}{2}]$，则 $a_1=0$；若 $x$ 属于靠右的（称为第二个，后同）子区间，则 $a_1=1$；
	
	\item 第二步：把 $x$ 所属的第一步中的小区间二等分为两半开闭子区间:  若 $x$ 属于第一个子区间，则 $a_2=0$；若 $x$ 属于第二个子区间，则 $a_2=1$；
	\item $\vdots$ 
	\item 第 $n$ 步： 把 $x$ 所属的上一步中的小区间二等分为两半开闭子区间:  若 $x$ 属于第一个子区间，则 $a_n=0$；若 $x$ 属于第二个子区间，则 $a_n=1$；
	\item $\vdots$
\ITEME
    注意到上述 $a_n$ 中，有无穷多个是 $1$ （为什么？）。这样我们就建立了 $(0,1]$ 和集合 $$S=\{(a_1, a_2,\cdots, a_n, \cdots): \mbox{每个 $a_n$ 是 $0$ 或 $1$ 且 其中无穷多个 $a_n$ 是 $1$} \}$$ 的一一对应关系。\\
2 决定 $\lvert S\rvert$：令 $$T=\{(a_1, a_2,\cdots, a_n, \cdots): \mbox{每个 $a_n$ 是 $0$ 或 $1$ 且 其中只有限多个 $a_n$ 是 $1$} \}$$ 和 $$ R=\{(a_1, a_2,\cdots, a_n, \cdots): \mbox{每个 $a_n$ 是 $0$ 或 $1$}\}$$ 
不难看出 $\lvert R\rvert=2^{\aleph_0}, \lvert T\rvert=\aleph_0$ （为什么？）且 $S\cup T= R$。由上面关于基数的基本定理 $2$，我们有 $$\lvert S\rvert=\lvert S\cup T\rvert =\lvert R\rvert=2^{\aleph_0}=\aleph_1$$

\subsubsection{连续统假设}
 
 \no {\bf 连续统假设} 不存在一个无穷集合，其基数严格介于自然数集的基数 $\aleph_0$ 与实数集的基数 $\aleph_1$（即连续统的基数）之间。 \\
 \no {\bf 广义连续统假设} 不存在一个无穷集合，其基数严格介于 $\aleph_n$ 与 $\aleph_{n+1}, n=0, 1,2, \cdots $ 之间。

在数学界广泛采用的ZFC集合公理系统下：
\ITEMS
\item Godel 证明了该假设的相容性：即不能证明它是错的； 
\item Cohen 证明了它的独立性：即不能证明它是对的。
\ITEME

\iffalse 
\section{$\mathbb{R}^n$上的拓扑}
这一部分的概念还不完整

\subsection{开集与闭集}

\subsubsection{开球}

我们把在$\mathbb{R}^n$空间上的开球定义为
\[
	B(x,r)=\{y\in\mathbb{R}^n:\lvert y-x\rvert<r\}
\]
我们可以把开球组成的集合视为一组基

\subsubsection{开闭集}
对于$A\subseteq\mathbb{R}^n$，如果
\[
	A = \bigcup_{\alpha\in I} B_\alpha(x_\alpha,r_\alpha)
\]
我们就称$A$是开集，而$A^c\subseteq\mathbb{R}^n$是闭集，用$U$和$F$分别表示开集和闭集；特别的，空集$\varnothing$和全集$\mathbb{R}^n$互补，它们既是开集也是闭集
\subsubsection{开闭集的运算}
我们称开集组成的集合为开集簇，闭集组成的集合为闭集簇
\[
	\{U_\alpha\subseteq\mathbb{R}^n\}_{\alpha\in I}\qquad\{F_\alpha\subseteq\mathbb{R}^n\}_{\alpha\in I}
\]

\ITEMS
\item $\bigcup\limits_{\alpha\in I} U_\alpha$仍是开集

\item $\bigcap\limits_{\alpha\in I} F_\alpha$仍是闭集

\item 开集的有限交集仍是开集

\item 闭集的有限并集仍是闭集
\ITEME

\subsubsection{内点与边界点}
对于$A\subseteq\mathbb{R}^n$，$x\in A$，存在包含于$A$的开球$B(x,r)\subseteq A$，使得$x\in B(x,r)$则我们称$x$是$A$的内点；如果对于任意包含$x$的开球$B(x,r)\ni x$，都有$B(x,r)$既包含$A$的点，也包含$A^c$的点，即$(B\cap A)\supset \{x\}$且$(B\cap A^c)\supset \{x\}$则称$x$是$A$的边界点
\fi

% 新章节



















\end{document}
















